\documentclass{article}
\usepackage[left=2cm, right=2cm, top=2cm, bottom=2cm]{geometry}
\usepackage{graphicx}
\usepackage{amsmath}
\usepackage{amssymb}
\usepackage{amsthm}
\usepackage{fancyhdr}
\usepackage{verbatim}
\usepackage{listings}
\usepackage{xcolor}
\usepackage{pdfpages}
\usepackage{cancel}
\usepackage{physics}
\usepackage{esint}

\lstset{
    language=C++,
    basicstyle=\ttfamily\footnotesize,
    keywordstyle=\color{blue},
    commentstyle=\color{green},
    stringstyle=\color{red},
    numbers=left,
    numberstyle=\tiny\color{gray},
    frame=single,
    breaklines=true,
    captionpos=b,
}


\title{MTH 553: Lecture \# 35}
\author{Cliff Sun}

\newtheorem{theorem}{Theorem}[section]
\newtheorem{lemma}[theorem]{Lemma}
\newtheorem{definition}[theorem]{Definition}
\newtheorem{conjecture}[theorem]{Conjecture}
\newtheorem{proposition}[theorem]{Proposition}
\newtheorem{corollary}[theorem]{Corollary}
\newtheorem{one minute paper}[theorem]{One Minute Paper}

\pagestyle{fancy}
\lhead{\textbf{\thepage}\ \ \nouppercase{\rightmark}}
\chead{MTH 553: Lecture \# 35}
\rhead{Cliff Sun}

\begin{document}

\maketitle

\section*{To sum it up}
We have a Poisson integral formula solving:
\begin{align}
    \Delta u = 0 & \Omega \\
    u = g & \partial \Omega
\end{align}

When $\Omega$ is the half space and when $\Omega$ is a ball. Note that $g$ is continuous and bounded. We take the normal derivative of the Greens function over $\Omega$. 

Using the formula for the greens function on the halfspace and ball, we can also solve Poisson's equation 

\begin{align}
    \Delta u = f & \Omega \\
    u = 0 & \partial \Omega
\end{align}

As 

\begin{equation}
    v(x) = \int_{\Omega}G(x,y)f(y)dy
\end{equation}

As for uniqueness, both the solutions of Poisson's and Laplace's equationa are unique. This can be proved using the max principle and energy methods. 

\section*{Harnack Inequality}

We will study qualitative properties of harmonic functions. Let $E \subset \Omega \subset \mathbb{R}^n$ be a compact subset. Then there exists $$C = C(\Omega, E)$$ such that for all 
harmonic functions $u$ on $\Omega$, if $u > 0$ in $\Omega$, then 

\begin{equation}
    \max_E u \leq C\min_E u \implies \frac{u(x)}{u(y)} \leq C
\end{equation}

i.e. positive harmonic functions have bounded oscillations. 
This actually isn't trivial, since this implies that the ratio is always constant (which is not the case for general bounded functions). 
This is a homework problem. \textbf{Hint:} use Poisson's formula for the ball. 

\begin{theorem}
    (Liouville's Theorem) If $u$ is harmonic on $\mathbb{R}^n$ and is bounded, then $u$ is constant.
\end{theorem}

\begin{proof}
    Let $a > 0$, then use the Poisson Integral formula for the ball $B(0,a)$ (origin with radius $a$)
    \begin{align}
        u(x) &= \frac{a^2 - |x|^2}{a \omega_n}\int_{\partial B(a)}\frac{u(y)}{|x-y|^n}dS(y)
    \end{align}
    Differentiate this with respect to $x_i$:
    \begin{align}
        u_{x_i}(x=0) &= 0 + \frac{a^2}{a \omega_n}\int_{\partial B(a)}\frac{ny_iu(y)}{|y|^{n+2}}\\
        |u_{x_i}(x=0)| &\leq \frac{a}{\omega_n}\frac{n}{|a^{n+1}|} ||u||_{L^{\infty}} (\omega_n a^{n-1})\\
        &= \frac{n}{a}||u||_{L^{\infty}}
    \end{align}
    If $a \rightarrow \infty$, then $u_{x_i}(x=0) \rightarrow 0$. Thus, $u$ is constant on $\mathbb{R}^n$. 
\end{proof}

\section*{Real Analyticity}

If $u$ is harmonic then it is real analytic (equals its Taylor Series) and hence in particular is $C^{\infty}$ smooth. 

\section*{Compactness}

Suppose $\{u_k\}$ is a uniform bounded sequence of harmonic functions in $\Omega \in \mathbb{R}^n$. Then if $\overline{B(x,r)} \subset \Omega$, then $\{u_k\}$ contains a 
subsequence which will converge uniformly on the ball and to some harmonic function. 

\begin{proof}
    Let $B(x,r)$ be in $\Omega$ and be a distance $a$ away from $\partial \Omega$. In other words, 
    \begin{equation}
        a = \text{dist}(\partial \Omega, \partial B(x,r))
    \end{equation}
    Therefore, $a > 0$. Let $M > 0$ such that $||u_k||_{L^{\infty}} \leq M$. Let $x_0 \in B(x,r)$ be the center of a ball $B(x,a)$ (Note, this ball DOESN'T hit the boundary). 
    By our argument from Liouville, we have that 
    \begin{equation}
        |\nabla u _k(x_0) | \leq \text{const.}\frac{nM}{a}
    \end{equation}
    Which is true for all $x_0 \in B(x,r)$. Now, by Arzela-Ascoli, we get that there exists a 
    uniformly convergent subsequence with continuous limit function. On homework 11, we will prove that this function is harmonic. (lol) Using the Poisson integral formula. 
\end{proof}

This result also says that the imbedding I: $\{\text{bounded harmonic functions on $\Omega$ with uniform norm}\} \rightarrow C(\overline{B})$ is compact. 

\end{document}